% !PRINTED 47
\weberpage{47}

We shall express the result obtained also as a theorem:

5. \gesperrt{The characters of a group $S$ can be combined into a group isomorphic with $S$.}

Let us also mention the theorem which follows immediately from the definition of $\chi_B$ \gesperrt{through} the formula (13):
\weq{20}{\chi_B(A) = \chi_A(B).}

Finally the following theorem also holds:

6. \gesperrt{If $A$ runs through all the elements of the group $S$, then for a fixed character $\chi$:}
\weq{21}{\overset{A}{\Sigma}\,\chi(A) = n \text{ or } = 0,}
\gesperrt{according as $\chi$ is the unit character or not. Likewise, if the element $A$ is held fixed and $\chi$ runs through the sequence of characters:}
\weq{22}{\overset{\chi}{\Sigma}\,\chi(A) = n \text{ or } = 0,}
\gesperrt{according as $A$ is the unit element or not.}

For the case that $\chi$ or $A$ are the unit elements of their groups, the formulae (21) and (22) are evident, since then each of the $n$ terms of the sum has the value 1. But if $\chi$ is not the unit character, then there is an element $B$ in $S$ such that $\chi(B)$ is not $= 1$. If we then set
\[
\overset{A}{\Sigma}\,\chi(A) = \sigma,
\]
it follows by multiplication with $\chi(B)$ that:
\[
\overset{A}{\Sigma}\,\chi(AB) = \sigma\chi(B).
\]
But since $AB$ runs through the whole group $S$ together with $A$, we have also $\overset{A}{\Sigma}\,\chi(AB) = \sigma$, hence $\sigma[1 - \chi(B)] = 0$ or $\sigma = 0$, q.\,e.\,d.

In the same way one proves the formula (22), which moreover also follows immediately from (21) by (20).
